\exercisesection{\S~1}

\begin{enumerate}
\def\labelenumi{\arabic{enumi})}
\item
  假设 \(k\) 的特征为 \(p > 0\) 。设 \(V\) 为向量空间，\(S\) 为有限集。映射 \(r: S \to \operatorname{End}(V)\) 满足拟交换条件 (AC) 当且仅当存在幂 \(q\) （\(p\) 的幂），使得 \([s^q, s'^q] = 0\) 对一切 \(s, s' \in S\) 成立。（利用第一章 §1 习题 19 公式 (1)。）
\item
  假设 \(k\) 是完全域。设 \(V\) 是有限维向量空间，\(u, v \in \operatorname{End}(V)\) 。设 \(u_s, u_n, v_s, v_n\) 是 \(u, v\) 的半单分量与幂零分量。下列条件等价：(i) 存在整数 \(m\) 使 \((\operatorname{ad} u)^m v = 0\) ；(ii) \(u_s\) 与 \(v\) 交换。（为证 (i) \(\Longrightarrow\) (ii)，化归到 \(k\) 代数闭的情形并利用引理 1 (ii)。）
\item
  我们采用第 2 节中的假设。假设 \(k\) 无限且拟交换条件 (AC) 成立。设 \(k'\) 是 \(k\) 的完全扩张。设 \(\lambda : S \to k\) 满足 \(V^{\lambda}(S) \neq 0\) 。置
\end{enumerate}

\[
\mathrm{V} ^ {\prime} = \mathrm{V} \otimes_ {k} k ^ {\prime}, \quad \mathrm{S} ^ {\prime} = \mathrm{S} \otimes_ {k} k ^ {\prime}.
\]

设 \(r' : S' \to \text{End}(V')\) 是由 \(r\) 经纯量扩张得到的线性映射。存在唯一的映射 \(\lambda' : S' \to k'\) 使得 \(V^{\lambda}(S) \otimes_k k' = V'^{\lambda'}(S')\) 。（化归到 \(V = V^{\lambda}(S)\) 的情形。设 \(P\) 是 \(S\) 上的多项式函数，\(q\) 是 \(k\) 的特征指数的一个幂且整除 \(\dim V\) ，使得 \(\lambda^q = P\) 。设 \(P'\) 是 \(S'\) 上扩张 \(P\) 的多项式函数。对每个 \(s' \in S'\) ，存在 \(\lambda'(s') \in k'\) 使得 \(\lambda'(s')^q = P'(s')\) 。证明 \(r'(s')\) 的特征多项式为 \((X - \lambda'(s'))^{\dim V}\) 。）

\begin{enumerate}
\def\labelenumi{\arabic{enumi})}
\setcounter{enumi}{3}
\tightlist
\item
  假设 k 的特征为零。设 \(\mathfrak{g} = \mathfrak{sl}(3, k)\) ，并用 a 表示由特征值为 1、-1、0 的对角矩阵生成的 g 的子代数。证明 a 在 g 中是约化的，a 在 g 中的交换子 m 由迹为零的对角矩阵组成，且 m 在 g 中的交换子等于 m，从而异于 a（参见第 5 节的注）。
\end{enumerate}

¶ 5) 假设 k 无限。设 g 为李代数，V 为 g-模。若 n 是整数 \(\geq 0\) ，用 \(V_{n}\) 表示由 \(v \in V\) 中使 \(x^{n}v = 0\) 对一切 \(x \in g\) 成立者所成的集。

\begin{enumerate}
\def\labelenumi{\alph{enumi})}
\tightlist
\item
  证明：若 \(v \in V_n\) ，\(x, y \in \mathfrak{g}\) ，则
\end{enumerate}

\[
\left(\sum_ {i = 1} ^ {n} x ^ {n - i} y x ^ {i - 1}\right) v = 0.
\]

（利用 \((x + ty)^n v = 0\) 对一切 \(t\in k\) 成立这一事实。）

在此公式中把 \(y\) 换成 \([x, y]\) ，据此\(^3\)推出 \((x^n y - y x^n)v = 0\) ，从而 \(x^n y v = 0\) 。

\begin{enumerate}
\def\labelenumi{\alph{enumi})}
\setcounter{enumi}{1}
\item
  证明 \(V_{n}\) 是 V 的 g-子模（利用 a)）。特别地，\(\mathrm{V}^{0}(\mathfrak{g})=\bigcup_{n}\mathrm{V}_{n}\) 是 V 的 g-子模。
\item
  假设 \(k = \mathbf{R}\) 或 \(\mathbf{C}\) ，用 \(G\) 表示以 \(\mathfrak{g}\) 为李代数的单连通李群；\(\mathfrak{g}\) 在 \(V\) 上的作用定义了 \(G\) 在 \(V\) 上的作用律（第三章 §6 第 1 节）。证明元素 \(v \in V\) 属于 \(V_n\) 当且仅当 \((s - 1)^n v = 0\) 对一切 \(s \in G\) 成立；特别地，\(V^0(\mathfrak{g}) = V^1(G)\) 。
\end{enumerate}

\begin{enumerate}
\def\labelenumi{\arabic{enumi})}
\setcounter{enumi}{5}
\tightlist
\item
  记号取自第一章 §3 习题 12。特别地，g 是李代数，M 是 g-模，而 \(\mathrm{H}^{p}(\mathfrak{g},\mathrm{M})=\mathrm{Z}^{p}(\mathfrak{g},\mathrm{M})/\mathrm{B}^{p}(\mathfrak{g},\mathrm{M})\) 是 g 的以 M 为系数的 p 次上同调空间。
\end{enumerate}

\begin{enumerate}
\def\labelenumi{\alph{enumi})}
\item
  证明 \(\mathrm{B}^{p}(\mathfrak{g},\mathrm{M})\) 与 \(\mathrm{Z}^{p}(\mathfrak{g},\mathrm{M})\) 在 g 的自然表示 \(\theta\) （作用于上链空间 \(\mathrm{C}^{p}(\mathfrak{g},\mathrm{M})\) ）之下稳定。由此得到 g 在 \(\mathrm{H}^{p}(\mathfrak{g},\mathrm{M})\) 上的一个表示。证明此表示是平凡的（利用公式 \(\theta = di + id\) ，见同前）。
\item
  设 \(\overline{k}\) 是 k 的代数闭包。设 \(x \in g\) ，并用 \(x_{M}\) 表示 M 上相应的自同态。设 \(\lambda_{1}, \ldots, \lambda_{n}\) （相应地 \(\mu_{1}, \ldots, \mu_{m}\) ）是 \(\overline{k}\) 中 \(ad_{g}x\) （相应地 \(x_{M}\) ）的特征值，且按重数重复。证明自同态 \(\theta(x)\) 在 \(C^{p}(g, M)\) 上的特征值是诸 \(\mu_{j} - (\lambda_{i_{1}} + \cdots + \lambda_{i_{p}})\) ，其中 \(1 \leq j \leq m\) 且
\end{enumerate}

\[
1 \leq i _ {1} <   i _ {2} <   \dots <   i _ {p} \leq n.
\]

试利用 a) 推出：\(\mathrm{H}^{p}(\mathfrak{g},\mathrm{M})=0\) ，只要诸 \(\mu_{j}-(\lambda_{i_{1}}+\cdots+\lambda_{i_{p}})\) 均不为零。

\begin{enumerate}
\def\labelenumi{\alph{enumi})}
\setcounter{enumi}{2}
\tightlist
\item
  假设表示 \(\mathfrak{g} \to \operatorname{End}(\mathrm{M})\) 是忠实的，且 \(x\) 满足条件：
\end{enumerate}

\[
\mu_ {j _ {1}} + \dots + \mu_ {j _ {p}} \neq \mu_ {k _ {1}} + \dots + \mu_ {k _ {p + 1}}\tag{($S_{p}$)}
\]

对所有 \(j_{1},\ldots,j_{p},k_{1},\ldots,k_{p+1}\in[1,m]\) 。

证明此时有 \(\mathrm{H}^p (\mathfrak{g},\mathrm{M}) = 0\) （注意 \(\lambda_{i}\) ，即 \(\mathrm{ad}_{\mathfrak{g}}x\) 的特征值，形如 \(\mu_j - \mu_k\) ，并应用 \(b\) ）。

\begin{enumerate}
\def\labelenumi{\arabic{enumi})}
\setcounter{enumi}{6}
\tightlist
\item
  设 \(\mathfrak{g}\) 是幂零李代数，V 是 \(\mathfrak{g}\) -模且 \(V^0(\mathfrak{g}) = 0\) 。证明 \(H^p(\mathfrak{g}, V) = 0\) 对一切 \(p \geq 0\) 成立。（化归到 \(V = V^\lambda(\mathfrak{g})\) 且 \(\lambda \neq 0\) 的情形，并选取元素 \(x \in \mathfrak{g}\) 使 \(\lambda(x) \neq 0\) 。应用习题 6 b)，注意诸 \(\lambda_i\) 全为零而诸 \(\mu_j\) 都等于 \(\lambda(x)\) 。）
\end{enumerate}

重新得到命题 9 的推论（取 p = 1）。\(^{4}\)

¶ 8) 假设 k 的特征为 p \textgreater{} 0。设 g 是 k 上的李代数，具有基 \((e_{1},\ldots,e_{n})\) 。用 U 表示 g 的包络代数，C 表示 U 的中心。对 \(i = 1,\ldots,n\) ，选取一个非零 p-多项式 \(f_{i}\) ，其次数为 \(d_{i}\) ，使得 \(f_{i}(\mathrm{ad}e_{i}) = 0\) ；于是 \(f_{i}(e_{i}) \in \mathrm{C}\) ，参见第一章 §7 习题 5。置 \(z_{i} = f_{i}(e_{i})\) 。

\begin{enumerate}
\def\labelenumi{\alph{enumi})}
\item
  证明 \(z_{1},\ldots,z_{n}\) 代数无关。若 \(A=k[z_{1},\ldots,z_{n}]\) ，证明 U 是自由 A-模，以单项式 \(e_{1}^{\alpha_{1}}\ldots e_{n}^{\alpha_{n}}\) （其中 \(0\leq\alpha_{i}\leq d_{i}\) ）为基。（利用 Poincaré-Birkhoff-Witt 定理。）U 在 A 上的秩 {[}U:A{]} 等于 \(d_{1}\ldots d_{n}\) ；它是 p 的幂。由此推出 C 是有限型 A-模，从而是有限型 k-代数且维数为 n（《交换代数》第八章）。
\item
  设 \(\mathrm{K}\) 是 \(\mathrm{A}\) 的分式域，并置
\end{enumerate}

\[
\mathrm{U} _ {(\mathrm{K})} = \mathrm{U} \otimes_ {\mathrm{A}} \mathrm{K}, \quad \mathrm{C} _ {(\mathrm{K})} = \mathrm{C} \otimes_ {\mathrm{A}} \mathrm{K}.
\]

则有 \(\mathrm{U}_{(\mathrm{K})}\supset\mathrm{C}_{(\mathrm{K})}\supset\mathrm{K}\) 。证明 \(\mathrm{U}_{(\mathrm{K})}\) 是一个域，其中心为 \(\mathrm{C}_{(\mathrm{K})}\) ，且它是 U 的（左、右）商域，参见第一章 §2 习题 10。推出 \([\mathrm{U}_{(\mathrm{K})}:\mathrm{C}_{(\mathrm{K})}]\) 具有 \(q^{2}\) 的形式，其中 q 是 p 的幂；我们有 \([\mathrm{C}_{(\mathrm{K})}:\mathrm{K}]=q_{\mathrm{C}}\) ，其中 \(q_{C}\) 是 p 的幂，而 \([U:A]=q_{C}q^{2}\) 。

\begin{enumerate}
\def\labelenumi{\alph{enumi})}
\setcounter{enumi}{2}
\tightlist
\item
  设 d 是 A 的非零元素，\(\Lambda\) 是 \(\mathrm{U}_{(\mathrm{K})}\) 的子环且 \(U \subset \Lambda \subset d^{-1}U\) 。证明 \(\Lambda = U\) 。{[}若 x = b/a（\(a \in A - \{0\}\) ）是 \(\Lambda\) 的元素，对 m 作归纳证明：关系 \(b \in Ua + U_m\) 蕴含 \(b \in Ua + U_{m-1}\) ，其中 \(\{U_m\}\) 是 U 的典则滤链。（为此，利用 gr U 整闭这一事实，并按《交换代数》第五章 §1 第 4 节 命题 15 的方式论证。）对 m = 0，即得 \(b \in Ua\) ，亦即 \(x \in U\) 。{]}
\end{enumerate}

推出 C 是整闭的。

\begin{enumerate}
\def\labelenumi{\alph{enumi})}
\setcounter{enumi}{3}
\tightlist
\item
  假设 k 代数闭。设 \(\rho: \mathfrak{g} \to \mathfrak{gl}(V)\) 是 g 的不可约线性表示，\(\rho_{U}\) 是 U 的相应表示。\(\rho_{U}\) 在 C 上的限制是从 C 到 k 的同态 \(\gamma_{\rho}\) （k 与 V 的位似群等同）；用 \(\alpha_{\rho}\) 表示它在 A 上的限制。证明：对每个从 k-代数 A （相应地 C）到 k 的同态 \(\alpha\) （相应地 \(\gamma\) ），至少存在 g 的一个不可约表示 \(\rho\) 使 \(\alpha_{\rho} = \alpha\) （相应地 \(\gamma_{\rho} = \gamma\) ），且这样的表示（在等价意义下）只有有限多个。证明 \(\dim V \leq q\) ，记号同 b)。\(^{5}\)
\end{enumerate}

¶ 9) 沿用上一习题的记号，并进一步假设 g 是幂零的。

\begin{enumerate}
\def\labelenumi{\alph{enumi})}
\item
  证明可以选取基 \((e_{1},\ldots,e_{n})\) ，使得对任意一对 \((i,j)\) ，\([e_{i},e_{j}]\) 是诸 \(e_{h}\) （ \(h > \sup(i,j)\) ）的线性组合。以下假设诸 \(e_{i}\) 满足此条件。对 \(i=1,\ldots,n\) 选取 p 的幂 \(q(i)\) 使 \(\operatorname{ad}(e_{i})^{q(i)} = 0\) ，并置 \(z_{i} = e_{i}^{q(i)}, A = k[z_{1},\ldots,z_{n}]\) ，参见习题 8. b)。设 \(\rho : g \to gl(V)\) 是 g 的线性表示。假设 \(\rho(e_{i})\) 对 \(i = 1,\ldots,n\) 幂零。证明 \(\rho(x)\) 对一切 \(x \in g\) 幂零。（对 \(n = \dim g\) 作归纳论证，并化归到 \(\rho\) 不可约的情形。证明在此情形 \(\rho(e_{n}) = 0\) ，然后应用归纳假设。）
\item
  设 \(\rho_1: \mathfrak{g} \to \mathfrak{gl}(V_1)\) 与 \(\rho_2: \mathfrak{g} \to \mathfrak{gl}(V_2)\) 是 \(\mathfrak{g}\) 的两个线性表示。假设 \(V_1\) 与 \(V_2\) 均 \(\neq 0\) ，且 \(V_1 = V^{\lambda_1}(\mathfrak{g}), V_2 = V^{\lambda_2}(\mathfrak{g})\) ，其中 \(\lambda_1\) 与 \(\lambda_2\) 是 \(\mathfrak{g}\) 上的两个函数，参见第 3 节。证明：若 \(\lambda_1(e_i) = \lambda_2(e_i)\) 对 \(i = 1, \ldots, n\) 成立，则 \(\lambda_1 = \lambda_2\) ，且存在一个非零 \(\mathfrak{g}\) -同态，从 \(V_1\) 到 \(V_2\) （把 \(b\) 应用于 \(\mathfrak{g}\) -模 \(V = \mathcal{L}(V_1, V_2)\) ，并用 Engel 定理证明 \(V\) 含有非零的 \(\mathfrak{g}\) -不变元）。推出：若此外 \(V_1\) 与 \(V_2\) 还是单模，则它们同构。
\item
  假设 k 代数闭。设 R 是 g 的不可约表示的等价类全体所成的集。若 \(\rho \in R\) ，置
\end{enumerate}

\[
x _ {\rho} = (x _ {\rho} (1), \dots , x _ {\rho} (n)) \in k ^ {n},
\]

其中 \(x_{\rho}(i)\) 是 \(\rho(e_{i})\) 的唯一特征值。证明 \(\rho \mapsto x_{\rho}\) 是从 R 到 \(k^{n}\) 的双射。（单性由 c) 得出，满性由习题 8 d) 得出。）由此推出下列结论：

\begin{enumerate}
\def\labelenumi{(\roman{enumi})}
\item
  对 A 的任一极大理想 \(\mathfrak{m}\) ，U/mU 模其根基的商是矩阵代数。
\item
  \(\mathfrak{g}\) 的任一不可约表示的次数都是 \(p\) 的幂（这由 (i) 及 \([\mathrm{U} / \mathfrak{mU}:k]\) 是 \(p\) 的幂这一事实得出）。
\item
  每个从 A 到 k 的同态都可唯一扩张为从 C 到 k 的同态（利用 C/mC 包含于 U/mU 的中心之中这一事实，而该中心是剩余域为 k 的局部 k-代数）。
\item
  存在整数 N ≥ 0，使 \(x^{p^{N}} \in A\) 对一切 \(x \in C\) 成立（这由 (iii) 得出）。\(^{6}\)
\end{enumerate}

¶ 10) 假设 \(k\) 的特征为 \(p > 0\) 。用 \(\mathfrak{g}\) 表示以 \(\{e_1, e_2, e_3\}\) 为基的李代数，其中 \([e_1, e_2] = e_3\) ，\([e_1, e_3] = [e_2, e_3] = 0\) 。

\begin{enumerate}
\def\labelenumi{\alph{enumi})}
\item
  证明 \(\mathrm{U}\mathfrak{g}\) 的中心是 \(k[e_1^p,e_2^p,e_3]\) 。
\item
  假设 k 代数闭。证明：对一切 \((\lambda_{1},\lambda_{2},\lambda_{3})\in k^{3}\) ，存在（在等价意义下）唯一的 g 的不可约表示 \(\rho\) ，使 \(\lambda_{i}\) 是 \(\rho(e_i)\) 的唯一特征值（ \(i = 1,2,3\) ）；\(\rho\) 的次数为 \(p\) ，若 \(\lambda_3 \neq 0\) ；为 1，若 \(\lambda_3 = 0\) 。（应用习题 8 与 9，或直接论证。）
\end{enumerate}

\begin{enumerate}
\def\labelenumi{\arabic{enumi})}
\setcounter{enumi}{10}
\tightlist
\item
  设 \(\mathfrak{h}\) 为幂零李代数，\(V\) 为不缩为 0 的 \(\mathfrak{h}\) -模，\(\lambda\) 为 \(\mathfrak{h}\) 上的函数且 \(V = V^{\lambda}(\mathfrak{h})\) 。证明下列性质等价：
\end{enumerate}

\begin{enumerate}
\def\labelenumi{(\roman{enumi})}
\item
  \(\lambda\) 是 \(\mathfrak{h}\) 上的线性形式，在 \([\mathfrak{h},\mathfrak{h}]\) 上为零。
\item
  存在 \(\mathrm{V}\) 的一组基，使得由 \(\mathfrak{h}\) 的元素所定义的自同态关于这组基都是三角形的。
\end{enumerate}

（为证 (i) \(\Longrightarrow\) (ii)，对 \(\mathfrak{h}\) -模 \(\mathcal{L}(\mathrm{W}_{\lambda},\mathrm{V})\) 应用 Engel 定理，其中 \(\mathrm{W}\) 是 1 维 \(\mathfrak{h}\) -模，由 \(\lambda\) 所定义。）

当 k 的特征为 0 时，性质 (i) 与 (ii) 成立（命题 9）。

